求完整分布时，数清等可能的排列、子集或间隔配置；只求均值时，常可把总数拆成指示变量之和，无须独立。方差还需两事件的联合概率，可按对象是否重叠分类。圆环计数要说明标记方法及每个对象被重复计数的次数。

交换性不等于独立；泊松总次数带来的独立分解也须由联合分布证明。项数依赖观测时，应保留这种依赖，用首步分解求随机和的矩。

高维体积可写成独立均匀样本满足约束的概率。大数定律定位阈值；在均值抵消且方差有限非零时，再用中心极限定理分析阈值附近的波动。
\par\smallskip
